\subsection{Resilience Against Deliberate State Compromise}
\label{sec:6.4.4}
A government does not need to out-coordinate a set of cooperating systems, fool a model with an adversarial example, or touch its weights, when it can compel the people who hold the model.

That mechanism has already been tested in a jurisdiction with about as much protection against it as exists anywhere. In February 2016, a federal magistrate in California ordered Apple, under the All Writs Act of 1789, to write and digitally sign new software defeating the passcode protections on an iPhone recovered from one of the San Bernardino shooters — code Apple had never built and did not want to build \autocite{pym2016applefbi}. Apple refused and moved to vacate the order, arguing in part that the authority claimed would not stop at one phone. The question was never decided: the FBI found a third party able to open the phone without Apple's help, and the government withdrew the request before any court ruled on whether the compulsion was lawful. Whether a democracy with an independent judiciary can force an unwilling company to compromise a system it built remains legally untested; a government less constrained by that judiciary does not need to win the argument first.

Section~\ref{sec:10.4} documents China's 2023 rules on generative AI providers — the same leverage Apple's lawyers argued no court should grant, exercised not against a locked phone but against the models themselves, as a standing condition of operating in that jurisdiction.

Section~\ref{sec:9.2} makes the parallel point about data, where a federated design raises the price of seizure by leaving no central store to seize — a protection that section~\ref{sec:6.4.2} finds inconveniencing rather than defeating. The same logic extends, partially, to a trained model. Weights already published and distributed beyond a developer's control cannot be recalled by a court order aimed at that developer, whatever the order compels going forward, and maintainers spread across separate jurisdictions raise the cost of coercion the same way — a government can compel the entity within its reach, not the ones outside it.

None of that removes the leverage; it only raises its price. A developer with employees, offices, and legal existence inside a government's borders answers to that government's courts regardless of how the model itself is built, and no architectural choice changes who can be held in contempt, arrested, or shut down for refusing an order. The Apple case tested that leverage against a single locked phone and left the question unresolved. Nobody has tested it yet against a model built specifically to resist the government making the request.
